\scanpage{59}
\deleted{
% E: text|upper left first paragraph|Par une petite modification
Par une petite modification de ce parcours, on montre alors que non seulement $\pi_1(X)\longrightarrow\pm1$, mais aussi $\pi_1(\widetilde X)\longrightarrow\pm1$ est surjectif; \unsure{je dis|je dirai} dans le cas $n$ pair, cela implique que $\pi_1(\unsure{\widetilde X|\widetilde{\widetilde X}})\longrightarrow\mathbb D_{2n}$ est surjectif (mais peut-être pas $\pi_1(\widetilde X)\longrightarrow\mathbb D_2$ ?)
% E: text|left middle heading|Calcul de la parité
Calcul de la parité du nb d'opérations élémentaires pour passer de $(D_i^s,u)$ à $(D_i^s,-u)$.
% E: text|middle introductory line|Le nb d'opérations
Le nb d'opérations est égal à :
% E: display|middle three lines|2(nb de sommets à balayer)
\[
\begin{gathered}
2(\text{nb de sommets à balayer})\\
+2(\text{nb de droites par dessus lesquelles}\ \text{balaye.})\\
+1\ (\text{décollage})+1\ (\text{atterrissage})
\end{gathered}
\]
% E: drawing|left middle beside count|horizontal line point s upward arrow
\figure{\begin{tikzpicture}[scale=.8]
\draw (-1.5,0)--(-.18,0) (.2,0)--(2.9,0);\fill (0,0) circle(.06);\node[below] at (0,0) {$s$};\draw[->] (.8,.2)--(.7,.85);
\end{tikzpicture}}
% E: margin note|middle right attached by brace to count|NB chaque droite
\margin{NB chaque droite par dessus laquelle on balaye compte comme un \unsure{demi} sommet par lequel on balayerait}
% E: text|lower left exclamation|merveilleux
merveilleux, c'est toujours pair !
% E: text|lower left cancelled opening|En fait par transport parallèle
\deleted{En fait, par transport parallèle, on trouve \unsure{bien ce pt.}}
% E: text|bottom left qualification|A vrai dire
A vrai dire, le raisonnement \unsure{est conclusif|n'est conclusif} parce que il faut \unsure{se placer ici} dans le cas des $\unsure{L|X}$ qui n'ont pas de faces supersingulières (pas des ``bandes supersingulières'') car celles-ci sont d'indice $2(n-1)$, non $2n$. Donc il vaut mieux
% E: text|right top continuation|partir d'une portion
partir d'une position proche de $D_i$ passant par $s$, par un petit pivotement autour de $s$.
% E: text|right upper middle revised count|Maintenant toute droite
Maintenant, toute droite à franchir compte comme \uline{une seule} opération; \unsure{du lg|de là}, et chaque droite \inserted{passant par $s$} \deleted{franchie} \inserted{est} franchie une fois et une seule. Donc le nb correct d'opérations est
% E: display|right middle two lines|2(nb de sommets non liés en s)+ordre de s
\[
\begin{gathered}
2(\text{nb de sommets non}\ \uline{\text{\unsure{ligués|liés}}}\ \text{en }s)\\
+\text{ordre de }s.
\end{gathered}
\]
% E: text|right lower conclusion|donc l'image du circuit
 donc l'image du circuit dans $\pm1$ est $-1$ ou $+1$, suivant que l'ordre \deleted{\unsure{de}} $\nu(s)$ de $s$ est pair ou impair.
Si donc il y a au moins un sommet simple, on est \unsure{sûr que l'on a bien} $\pi_1(\unsure{\widetilde X|\widetilde{\widetilde X}})\longrightarrow\pm1$ surjectif, et encore $\pi_1(\widetilde X)\longrightarrow\unsure{\pm1|\mathbb D_{2n}}$ surj.
}
\mspage{58}
